% Einheit 4 von 4 – Ungleichungen (bank/gleichungen-loesen/e4.jsonl, zusammenbau v0.1)
%% TODO Zeitmarke Einheit 4: Marken-Zeile nicht lesbar
%% TODO Zweigzeile Teil 1: was hier gelernt wird (Satz in Schülersprache aus dem Einheitstitel) – Einheit 4
\einheitenkopf[e4]{Einheit 4 von 4 · Ungleichungen}
\zweigzeile{Abitur GK}
\setcounter{aufgabe}{19}

%% TODO Titel: Ich-kann-Satz fehlt in der Bank (Platzhalter: Kettenname) – Nr. 20
%% TODO Anweisung: ein Satz über der ersten Teilaufgabe fehlt in der Bank – Nr. 20
\begin{aufgabe}{Ungleichungen}
%% TODO \rechenplatz{n}: Schrittzahl des Grundfalls fehlt in der Bank (nur bei fester Schreibform, 2.3 b)
\begin{teile}
\teil Welche Lösungen zählen? $x^2 \cdot (x - 3) > 0$ \kreuz{$x > 3$} \\ \kreuz{$x > 3$ und $x = 0$} \\ \kreuz{$x \ge 3$}
\teil Für welche $x$ gilt $(x - 1) \cdot (x + 2)^2 > 0$? \leerfeld
\end{teile}
\begin{gleichungsraster}[2]
\gl{\text{$f(x) = x^3 - 6x^2 + 11x - 1$, $g(x) = 0{,}5x^3 - 2x^2 + 3x - 1$. Bestimme rechnerisch alle $x$ mit $f(x) > g(x)$. (Abitur 2023 GK)}} &
\gl{\text{Die Düngerkonzentration im Boden ist $w_k(t) = k \cdot t \cdot e^{-0{,}1t} + 20$ ($t$ in Tagen, $k > 0$). Ab welchem $k$ beträgt sie am 15. Tag mindestens 40? (Abitur 2019 GK)}} \\
\end{gleichungsraster}
\begin{teile}
\teil $f(x) = 0{,}25x^4 - 2x^2 + 4$ wird für $x \ge 0$ durch $p(x) = -2x^2 + 4$ angenähert. Die Näherung ist brauchbar, solange der Betrag der Differenz höchstens $0{,}25$ ist. Beschreibe einen Lösungsweg, mit dem man alle $x \ge 0$ findet, für die $p$ brauchbar ist. \hfill (Abitur 2024 GK)
\end{teile}
\end{aufgabe}

%% TODO Titel: Ich-kann-Satz fehlt in der Bank (Platzhalter: Kettenname) – Nr. 21
%% TODO Anweisung: ein Satz über der ersten Teilaufgabe fehlt in der Bank – Nr. 21
\begin{aufgabe}{Intervall, in dem eine Modellfunktion mindestens einen vorgegebenen Wert annimmt, über eine Ungleichung bestimmen}
\begin{gleichungsraster}[2]
\gl{\text{Beim Erhitzen einer Flüssigkeit gilt $f(t) = 20 + 15t \cdot e^{-\frac{t}{8}}$ ($t$ in min, $f$ in °C). In welchem Zeitraum ist sie mindestens 60 °C heiß? Löse die Ungleichung mit dem Rechner (zwei Stellen). (Abitur 2017 GK)}} & \\
\end{gleichungsraster}
\end{aufgabe}

%% TODO Titel: Ich-kann-Satz fehlt in der Bank (Platzhalter: Kettenname) – Nr. 22
%% TODO Anweisung: ein Satz über der ersten Teilaufgabe fehlt in der Bank – Nr. 22
\begin{aufgabe}{Ungleichungen – Fehler finden, Begründen, Anwendung, Darstellungswechsel}
\begin{teile}
\teil Finde den Fehler und gib die richtige Lösungsmenge an. \rechnung{x \cdot (x - 6)^2 &> 0 \\ x &> 0}
\teil Begründe, warum der Faktor $(x - 7)^2$ das Vorzeichen von $(x + 2) \cdot (x - 7)^2$ nicht ändert, außer bei $x = 7$.
\end{teile}
\begin{gleichungsraster}[2]
\gl{\text{Der Gewinn einer Firma ist $G(x) = -x^2 + 12x - 20$ ($x$ in tausend Stück, $G$ in tausend €). Bei welchen Mengen macht sie Gewinn?}} &
\end{gleichungsraster}
\begin{teile}
\teil Die Abbildung zeigt die Graphen von $f$ und $g$. Lies ab, für welche $x$ gilt $f(x) > g(x)$.

\begin{ksys}[xmin=-3,xmax=5,ymin=-1,ymax=5,ablesen] \parabel{-0.5}{1}{4}{f} \gerade{0}{2}{g} \end{ksys}
\leerfeld
\end{teile}
\end{aufgabe}
